% ===========================================================================
%  sviluppatore/03-architettura/02-strangler.tex
% ===========================================================================
\section{Due applicazioni web, e il motivo}\label{sez:strangler}

Dentro la metà locale ci sono due applicazioni web (\vediFig{strangler}):

\begin{center}
\begin{tabular}{@{}L{34mm}L{24mm}L{80mm}@{}}
\toprule
 & \textbf{Flask} & \textbf{FastAPI} \\
\midrule
\rigaTenue
File principale & \file{dashboard/app.py} & \file{api/app.py} \\
Porta           & \num{5000}              & \num{8000} \\
\rigaTenue
Dimensione      & \num{249} rotte         & \num{13} router \\
Restituisce     & pagine HTML e JSON      & solo JSON \\
\rigaTenue
Autenticazione  & guardia globale, quasi tutto aperto & nessuna \\
Parla con       & direttamente il motore  & i casi d'uso, via
                                            \simbolo{Container} \\
\bottomrule
\end{tabular}
\end{center}

\subsection{Che cos'è uno strangler fig}

La pianta da cui prende il nome cresce attorno a un albero ospite, lo
avvolge e lo sostituisce gradualmente: alla fine resta in piedi da sola, e
l'albero di dentro non c'è più. Nessun momento in cui il bosco resta senza
albero.

Applicato qui: la dashboard Flask ha dentro il codice che è stato scritto per
primo, con il motore chiamato direttamente. L'API FastAPI è il nuovo guscio e
parla solo con i casi d'uso. Le due convivono, e il confine si sposta una
funzionalità alla volta.

\begin{perche}{Perché non si riscrive la dashboard e basta}
Perché la dashboard funziona, e perché \num{28494} righe di Python più
\num{42072} di template non si riscrivono con la garanzia di non perdere
niente. Una riscrittura in blocco ha due esiti noti: o dura abbastanza da
essere abbandonata a metà, o arriva in fondo avendo perso silenziosamente i
casi limite che nessuno aveva scritto da nessuna parte --- e in un sistema che
muove denaro, un caso limite perso è un ordine sbagliato.
\end{perche}

\subsection{Come si avvia l'API}

\begin{codiceInterno}{python}{Il montaggio dell'applicazione}{api/app.py}
def create_app(container: Optional[Container] = None, jobs=None) -> FastAPI:
    container = container or build_container()
    app = FastAPI(title="QuantSys Engine API", version="0.2.0")
    app.state.container = container

    if jobs is None:
        # JobManager reale: max_workers=1 = serializzato (default sicuro del progetto)
        from execution.job_manager import JobManager
        job_manager = JobManager(max_workers=1)
        app.state.job_manager = job_manager
        jobs = JobManagerJobService(container.run_walkforward, job_manager)
    app.state.jobs = jobs

    app.include_router(health.router)
    app.include_router(strategies.router)
    ...
\end{codiceInterno}

\noindent
Due cose meritano attenzione. La prima è che \simbolo{create\_app} accetta un
\simbolo{Container} dall'esterno: nei test si passa un contenitore di finti e
l'API intera gira senza toccare il disco. La seconda è
\cod{max\_workers=1}: i job sono \emph{serializzati}, ed è il default sicuro
del progetto. Due backtest in parallelo si contendono la stessa cache dati e
gli stessi file di output.

\subsection{I tredici router}

\begin{center}
\begin{tabular}{@{}L{32mm}L{112mm}@{}}
\toprule
\textbf{Router} & \textbf{Che cosa espone} \\
\midrule
\rigaTenue
\file{health}             & Stato del servizio. \\
\file{strategies}         & Elenco e dettaglio delle strategie. \\
\rigaTenue
\file{strategy\_admin}    & Validazione e installazione di una strategia. \\
\file{runs}               & Avvio di un walk-forward, stato, log in SSE. \\
\rigaTenue
\file{streams}            & Catalogo degli stream prodotti. \\
\file{ensemble}           & Composizione: pesi, overlay, metriche. \\
\rigaTenue
\file{models}             & Registro dei modelli ML e promozione. \\
\file{factor\_attribution}& Scomposizione del rendimento per fattore. \\
\rigaTenue
\file{stress\_test}       & Prove sotto scenario. \\
\file{live}               & Controllo dell'operatività. \\
\rigaTenue
\file{generator}          & Generazione di codice da descrizione. \\
\file{nodes}              & Compilazione del grafo a nodi in Python. \\
\rigaTenue
\file{fidelity}           & Confronto fra trade disegnati a mano e strategia. \\
\bottomrule
\end{tabular}
\end{center}

\begin{attenzione}{Le rotte FastAPI non esistono per Flask}
Sono due applicazioni separate, in due processi. Nessuna rotta di \pkgApi{} è
registrata nell'app Flask, e viceversa. Cercare \percorsoHTTP{/api/strategies} sulla
porta \num{5000} e trovarci una rotta Flask con lo stesso nome ma
comportamento diverso è possibile, ed è successo: quando si documenta o si
collauda un endpoint, la porta fa parte dell'indirizzo.
\end{attenzione}

\subsection{Dove passa il confine, oggi}

La regola pratica: \textbf{ciò che è nuovo nasce nel guscio}. Una funzionalità
nuova si scrive come caso d'uso, si espone dall'API, e la dashboard la chiama
come chiamerebbe qualunque servizio. Ciò che esiste si sposta quando lo si
deve toccare comunque, non prima e non per ordine --- spostare codice che
funziona e che nessuno sta modificando è rischio senza ritorno.
